\documentclass[text06.tex]{subfiles}
\begin{document}
%偶数ページに飛ばしたい
\newpage
\section{音声でLEDのON OFFをする}

認識した言葉でLEDを動かしてみましょう。「ON」「OFF」と話しかけると、LEDがついたり消えたりします。

\begin{lstlisting}[caption=onoff.yomi,label=onoff.yomi,aboveskip=0.4em,belowskip=0.3em]
ON おん
OFF おふ
\end{lstlisting}

\enlargethispage{3\baselineskip}
\begin{lstlisting}[caption=ledbyvoice.hsp,label=ledbyvoice.hsp,aboveskip=0.4em,belowskip=0.2em]
#include "hsp3dish.as"
#include "julius.as"

redraw 0
pos 0, 0
mes "Loading Julius..."
redraw 1

init_julius "onoff.dic"

sockidx = 0
domain = "127.0.0.1"
PORT = 10500

sockopen sockidx, domain, PORT

if stat != 0 {
  redraw 0
  mes "Error in opening socket"
  redraw 1
  wait 30
  stop
}

sdim words, 4096, 0
ddim cm, 0
repeat -1
  if is_recieved(sockidx) = 0 {
    get_word_list words, cm
    repeat length(words)
      if cm(cnt) > 0.7 {
        if words(cnt) = "ON" {			<#blue#;認識された文字がONのとき#>
          gpio 17, 1				<#blue#;GPIO17のLEDを光らせる#>
        }else : if words(cnt) = "OFF" {		<#blue#;認識された文字がOFFのとき#>
          gpio 17, 0				<#blue#;GPIO17を消す#>
        }
      }
    loop
  }
  redraw 0
  pos 0,0
  mes "話しかけてLEDのONとOFFをしよう"
  redraw 1
  await 30
loop
\end{lstlisting}

 \textasciitilde /07に\mbox{onoff.yomi}ファイルがあります。これは音声でLEDのON/OFFをするための辞書のもとになるファイルです。このファイルはリスト\ref{onoff.yomi}のように書いてあります。ONとOFFがコンピュータの使うラベル、おんとおふが読み方です。これを\mbox{\code{convert\_yomi.sh}}で変換したものが\mbox{onoff.dic}です。 \textasciitilde /07/\mbox{ledbyvoice.hsp} を開いてください。

このプログラムでおさえるのは、次の4つです。

\begin{itemize}
\item \code{words(cnt)} に、認識した言葉が入る
\item \code{"ON"} ならLEDをつける
\item \code{"OFF"} ならLEDを消す
\item \code{cm} が一定以上の結果だけ使う
\end{itemize}

\code{cm(cnt) > 0.7} の 0.7 は、これ以上なら使う、と決める\ruby{境目}{さかい|め}の値です。このような境目の値を「しきい値」といいます。しきい値を変えると、LEDの反応のしやすさが変わります。

\begin{tcolorbox}[title=\useOmetoi]
\begin{enumerate}
\addex{ターミナルを開いて、 \textasciitilde /07/に移動しましょう。}
\addex{\mbox{ledbyvoice.hsp}をHSPスクリプトエディタで実行しましょう。「ON」「OFF」と話しかけて、LEDがつく・消えることを確かめましょう。}
\addex{点灯させるLEDを変\ruby{更}{こう}しましょう。}
\addquiz{\mbox{ledbyvoice.hsp}の \code{cm(cnt) > 0.7} を、0.3、0.7、0.9 に変えて実行し、ちがいを\ruby{観察}{かん|さつ}しましょう。\\それぞれの値で「ON」「OFF」と、辞書にない言葉（例えば「こんにちは」）を話しかけてみましょう。\\値を小さくするとLEDは反応しやすくなりましたか。関係のない言葉でも反応することがありましたか。値を大きくするとどうなりましたか。}
\addex{時間があったらやりましょう。\\LEDを\ruby{操作}{そう|さ}する言葉を「つける」と「けす」に変更しましょう。\\ヒント：\mbox{onoff.yomi} の読み方を変え、\mbox{\code{convert\_yomi.sh}} で \mbox{onoff.dic} を作り直します。}
\end{enumerate}
\end{tcolorbox}
\end{document}
